% HD-S21. GPT-6 Astra Pro, /root/sixfolds, 21 November 2026.
% Inherited effort is not exposed. Original proof; arithmetic separately checked.
\section{An actual coupled graph construction and its trace kernel}

\paragraph{HD-S21.}
For every positive squarefree integer $d$, the construction below gives
perfect complexes $E$ on an explicit CM abelian eightfold, with no
negative self-Ext and with
\[
 \operatorname{ch}(E)=e^\theta+10080w_d,
 \qquad 0\ne w_d\in W_{\mathbb Q(\sqrt{-d})}.
\]
It includes nonsplit extension triangles. Every one of these complexes
has a nonzero kernel in the full degree-two semiregularity map.
Its stabilizer in $A\times\widehat A$ is trivial, so restriction to
invariants for a finite subgroup of that point stabilizer does not
remove this kernel. The proof makes no assertion about arbitrary
longer complexes or other coupled constructions.

\paragraph{1. The variety, the action, and actual graph cycles.}
Put $X=E_i^4$ with its product principal polarization $H$ and
$A=X\times X$. Write $(x_j,y_j,u_j,v_j)$ for the integral symplectic
one-forms on the two copies, for $1\le j\le4$, and put
\[
 \theta=dH_1+H_2,\qquad
 J_d=\begin{pmatrix}0&I\\-dI&0\end{pmatrix}\in\operatorname{End}(A).
\]
Thus $J_d^2=-dI$. The action $s\mapsto J_d$, $s^2=-d$, is
skew-adjoint for the positive product polarization $\theta$.
Its two eigenvalues have holomorphic multiplicity four.
On the first-copy $a_j,b_j$ cycles the associated hermitian form has
matrix $\left(\begin{smallmatrix}0&ds\\-ds&0\end{smallmatrix}\right)$;
the basis $a_j,b_j/(ds)$ has matrix $U$. This exhibits the rational
form as $U^4$.

For each integer $t$, let
\[
 B_t=\{(z,tz):z\in X\}\subset A,\qquad
 F(t)=\prod_{j=1}^4(u_j-tx_j)(v_j-ty_j).
\]
The surjective homomorphism $(z,z')\mapsto z'-tz$ has connected
kernel $B_t$, so $F(t)=[B_t]$ integrally. Both tangent bundles and
the normal bundle are trivial; Grothendieck--Riemann--Roch therefore
gives $\operatorname{ch}(\mathcal O_{B_t})=F(t)$.

Since $J_d^*x_j=u_j$ and $J_d^*u_j=-dx_j$, each
$u_j-sx_j$ and $v_j-sy_j$ has eigenvalue $-s$. Consequently
\[
 w_d=\frac{F(s)-F(-s)}{2s}
\]
is a rational integral Weil class. It is nonzero: the coefficient
of a monomial using exactly one of the $x_j,y_j$ is $-1$.
Each pair satisfies
\[
 (u-sx)(v-sy)(u+sx)(v+sy)=-4d\,xyuv,
\]
so $\int_A w_d^2=128d^3$. In addition, $\theta F(\pm s)=0$,
as follows by multiplying the corresponding pair by $dxy+uv$.

\paragraph{2. Integral interpolation.}
Define the four signed integers
\begin{align*}
 b_1&=14(d+4)(d+9)(d+16),\\
 b_2&=-14(d+1)(d+9)(d+16),\\
 b_3&=6(d+1)(d+4)(d+16),\\
 b_4&=-(d+1)(d+4)(d+9).
\end{align*}
The exact identities
\[
 2\sum_{j=1}^4 b_jj^{2r+1}=10080(-d)^r
 \quad(0\le r\le3)
\]
are Lagrange interpolation at $1,4,9,16$: the weights before
clearing denominators are
\[
 \frac{b_j}{10080}
 =\frac1{2j}\prod_{\ell\ne j}
       \frac{-d-\ell^2}{j^2-\ell^2}.
\]
The odd part of $F$ has degrees $1,3,5,7$, and hence
\[
 \sum_{j=1}^4b_j\bigl([B_j]-[B_{-j}]\bigr)=10080w_d.
\]
The accompanying exact script checks all four polynomial identities
and their nine class-number-one specializations.

\paragraph{3. The actual complexes.}
Let $L$ be a line bundle of class $\theta$, and put $m_j=|b_j|$.
Define sheaves
\begin{align*}
 P={}&L\oplus\mathcal O_{B_1}^{m_1}
       \oplus\mathcal O_{B_{-2}}^{m_2}
       \oplus\mathcal O_{B_3}^{m_3}
       \oplus\mathcal O_{B_{-4}}^{m_4},\\
 N={}&\mathcal O_{B_{-1}}^{m_1}
       \oplus\mathcal O_{B_2}^{m_2}
       \oplus\mathcal O_{B_{-3}}^{m_3}
       \oplus\mathcal O_{B_4}^{m_4}.
\end{align*}
For any $\delta\in\operatorname{Ext}^4(P,N)$ take its cone shifted
by $-1$, so that the defining triangle is
\begin{equation}\label{hd-s21-triangle}
 N[3]\longrightarrow E\xrightarrow{q}P
       \xrightarrow{\delta}N[4].
\end{equation}
Every nonzero $\delta$ gives a nonsplit triangle. Such classes exist:
for distinct slopes $a,b$, the graphs intersect transversely in
$|a-b|^8$ points. The Koszul resolution in four complementary normal
coordinates gives
\[
 \operatorname{Ext}^k(\mathcal O_{B_a},\mathcal O_{B_b})=0
 \ (k\ne4),\qquad
 \dim\operatorname{Ext}^4(\mathcal O_{B_a},\mathcal O_{B_b})
       =|a-b|^8.
\]
Also $\operatorname{Ext}^k(L,\mathcal O_{B_a})
=H^k(B_a,L^{-1}|_{B_a})$ is concentrated in degree four, by
ampleness and Serre duality on $B_a$; duality on $A$ gives the same
concentration in the reverse direction. Thus
\begin{equation}\label{hd-s21-cross}
 \operatorname{Ext}^k(P,N)=\operatorname{Ext}^k(N,P)=0
 \quad(k\ne4).
\end{equation}
The objects $P$ and $N[3]$ have no negative self-Ext or negative
cross-Ext. Applying the long exact Hom sequences of
\eqref{hd-s21-triangle} in each variable proves
$\operatorname{Ext}^{<0}(E,E)=0$.
All sheaves are perfect on the smooth variety $A$, hence so is $E$.
Its character is exactly the expression in the statement.

\paragraph{4. A canonical injection and its trace.}
Applying $\operatorname{Hom}(P,-)$ to
\eqref{hd-s21-triangle}, and using the vanishing of
$\operatorname{Ext}^5(P,N)$ and $\operatorname{Ext}^6(P,N)$,
gives an isomorphism
\[
 q[2]_*:\operatorname{Ext}^2(P,E)
          \xrightarrow{\ \sim\ }\operatorname{Ext}^2(P,P).
\]
For $f\in\operatorname{Ext}^2(P,P)$ write $\alpha_f:P\to E[2]$
for its unique lift and set $j(f)=\alpha_fq$.
This is an injection into $\operatorname{Ext}^2(E,E)$.
Indeed, the kernel of precomposition by $q$ consists of maps
$\beta\delta$, with $\beta:N[4]\to E[2]$. But
$q[2]\beta=0$ since $\operatorname{Hom}(N[4],P[2])=0$.
The displayed isomorphism then makes $\beta\delta=0$.

Write $\sigma_T$ for the full degree-two Atiyah trace of an object
$T$. Naturality of its Atiyah class and cyclicity of the trace give,
component by component,
\begin{align*}
 \operatorname{tr}_E\!\left(
   \alpha_fq\,\frac{\operatorname{At}_E^a}{a!}\right)
 &=\operatorname{tr}_P\!\left(
   q[2]\alpha_f\,\frac{\operatorname{At}_P^a}{a!}\right),\\
 \sigma_E(j(f))&=\sigma_P(f).
\end{align*}
The maps have degrees zero and two, so cyclic interchange introduces
no odd sign. These are actual maps of Ext groups; there is no appeal
to equality of dimensions of parameter spaces.

For a graph block $\mathcal O_{B}^{m}$, its Atiyah class is
$\operatorname{At}_{\mathcal O_B}\otimes I_m$. Thus
\[
 \sigma_P(\eta\otimes M)=0
 \quad\text{for }\eta\in\operatorname{Ext}^2(\mathcal O_B,\mathcal O_B),
 \quad\operatorname{tr}(M)=0.
\]
For example choose a nonzero $\eta\in H^2(\mathcal O_B)$ and
$M=\operatorname{diag}(1,-1,0,\ldots,0)$. This is nonzero and
$m_1\ge11900$, so it is available in the first positive block.
Its image under $j$ is a nonzero class in $\ker\sigma_E$.
The graph automorphism $(z,z')\mapsto(z,z'-tz)$ identifies the
self-Ext of $\mathcal O_{B_t}$ with the K\"unneth product for
$\mathcal O_X\boxtimes\mathcal O_0$, giving dimension
$\binom82=28$ in degree two. In particular the proof gives the
explicit lower bound
\[
 \dim\ker\sigma_E\ \ge\ 28\sum_{j=1}^4(m_j^2-1)>0.
\]
The same two Hom sequences also give
\[
 0\longrightarrow\operatorname{Ext}^2(P,P)
 \xrightarrow{j}\operatorname{Ext}^2(E,E)
 \longrightarrow\operatorname{Ext}^2(N,N)\longrightarrow0.
\]
Indeed, $\operatorname{Hom}(N[3],E[2])$ identifies with
$\operatorname{Ext}^2(N,N)$ because the neighboring negative Ext
groups between $N$ and $P$ vanish. Its connecting map projects to
zero in $\operatorname{Ext}^3(P,P)$, while
$\operatorname{Ext}^3(P,E)\to\operatorname{Ext}^3(P,P)$ is an
isomorphism by \eqref{hd-s21-cross}. Consequently
$\dim\operatorname{Ext}^2(E,E)=28(1+2\sum_jm_j^2)$.
For clarity, this conclusion uses the trace identity above, and
does not infer injectivity or its failure by comparing the degree-two
character-contraction rank with the dimension of $\operatorname{Ext}^2$.

\paragraph{5. The exact point stabilizer.}
The only nonzero cohomology sheaves of $E$ are
$\mathcal H^{-3}(E)=N$ and $\mathcal H^0(E)=P$.
An isomorphism $t_a^*E\otimes Q\cong E$, with
$a\in A$ and $Q\in\widehat A$, must preserve each graph component
of the support of $N$: their distinct tangent spaces prevent a
translation from permuting different slopes. Hence
$a\in B_{-1}\cap B_2\cap B_{-3}\cap B_4$.
Writing $a=(z,z')$, the first three equations give $3z=2z=0$,
so $z=z'=0$. The quotient of $P$ by its torsion subsheaf is $L$;
the remaining isomorphism therefore gives $L\otimes Q\cong L$,
and $Q=\mathcal O_A$. This proves the asserted trivial stabilizer.

\paragraph{6. Replacing every multiplicity by distinct translated cosets.}
There is an equally explicit version with no repeated sheaf summands.
Choose the infinite-order point
$v=(\sqrt2\bmod\mathbb Z[i],0,0,0)\in X$.
Replace each $\mathcal O_{B_t}^{m_j}$ in $P$ and $N$ by
\[
 \bigoplus_{\ell=0}^{m_j-1}
   \mathcal O_{B_t+(0,\ell v)}.
\]
These cosets of each fixed graph are pairwise disjoint. Different
slopes remain transverse, so \eqref{hd-s21-cross}, the cone,
negative-Ext vanishing and the canonical injection $j$ remain valid.
Translation acts trivially on the cohomology of $A$, so the character
is unchanged.

Identify the self-Ext groups of the cosets by their actual translation
isomorphisms. Naturality of the Atiyah trace under translation gives
the same map from all these identified groups to the cohomology of
$A$. For two cosets of the same slope, the element with diagonal
entries $\eta$ and $-\eta$ therefore has zero trace. It is nonzero
in $\operatorname{Ext}^2(P,P)$ and stays nonzero under $j$.
This proves
\[
 \dim\ker\sigma_E\ge28\sum_{j=1}^4(m_j-1)>0
\]
for every extension class in this translated construction, and
$\dim\operatorname{Ext}^2(E,E)=28(1+2\sum_jm_j)$.

The point stabilizer is still trivial. For any fixed slope the set
of cosets is parametrized by
$S_j=\{0,v,\ldots,(m_j-1)v\}\subset X$ under the quotient
$q_t(z,z')=z'-tz$. A translation $c$ satisfying $S_j+c=S_j$ is
torsion, since it permutes a finite set. It is also an integer
multiple of $v$, since it is the difference of two elements of
$S_j$. The infinite order of $v$ implies $c=0$. Preservation of
the four negative-slope families therefore forces
$q_t(a)=0$ for all four slopes. The previous $2z=3z=0$ argument
again gives $a=0$, and the rank-one quotient $L$ again gives
$Q=\mathcal O_A$.

Thus replacing multiplicities by distinct translated cosets removes
the matrix-valued kernel from repeated copies, but the explicitly
identified diagonal cancellation classes still lie in the kernel.
The proof here covers precisely the displayed two-layer cones.

\paragraph{Sources and proof status.}
The graph interpolation idea was suggested by the unrefereed
Bhattacharjee--Bhattacharya preprint,
\emph{Secant Objects for Weil Classes on Abelian Varieties},
Example1, printed p.35.
\par DOI: \texttt{10.20944/preprints202611.0672.v1}.
\par
No assertion of that preprint is a premise here. The graph class,
interpolation, actual cones, negative-Ext vanishing, trace injection,
kernel and stabilizer have all been derived above.
The trace conventions are those checked from
Buchweitz--Flenner, Adv. Math.217 (2008), Proposition6.2.1,
Corollary6.3.2 and Theorem6.4.2, and Perry's permitted
semiregularity theorem. The polynomial script is an arithmetic
certificate only. A full author reread is recorded separately;
an external whole-proof review is requested.
